\section{Introducción: por qué están en auge los agentes de desarrollo de software}
\subsection{Contexto de la industria y plataforma}
Graham Neubig abre presentándose como profesor de la CMU y científico en jefe de All Hands AI, y relata su acumulación doble entre el desarrollo de software y la inteligencia artificial. All Hands AI mantiene el framework de código abierto OpenHands (antes OpenDev), cuyo objetivo es empaquetar la capacidad de «escribir código» en un flujo lineal que un agente pueda ejecutar, de modo que los desarrolladores puedan concentrarse en los requerimientos y la revisión, y dejarle al agente las pruebas, la búsqueda y los parches repetitivos.

Señala que una plataforma abierta permite que la comunidad construya en conjunto herramientas de depuración y repositorios de ejemplo, extraiga entradas de issues reales, dé seguimiento a las operaciones del agente en repositorios públicos y retroalimente las mejoras al modelo. Este ciclo cerrado permite que la investigación llegue directamente a la práctica de ingeniería, en lugar de quedarse en los puntos de referencia de laboratorio.

\subsection{Empoderar al público general}
Menciona en particular el artículo de Mark Andreessen de 2011, ``Why Software Is Eating the World'', que señala que todas las industrias están siendo redefinidas por el software: del entretenimiento a la agricultura, la defensa y los servicios médicos, el software ha ido sustituyendo procesos físicos, y el ámbito del que se encargará el software en el futuro solo va a crecer. Esto también explica por qué es necesario que más personas cuenten con la capacidad de construir software rápidamente, en lugar de que el desarrollo de software siga monopolizado por un pequeño número de ingenieros profesionales.

Al presentar la motivación, Graham usa como analogía varios avances científicos recientes impulsados por sistemas de software, para subrayar que el progreso científico contemporáneo está profundamente entrelazado con el código. Para él, el valor del agente está en poder transferir la capacidad de ingeniería a un grupo más amplio de investigadores, gerentes de producto y equipos interdisciplinarios, de modo que «crear software» se vuelva una capacidad generalizada en lugar de un recurso escaso.

\begin{knowledgebox}{La predicción de Mark Andreessen}
El artículo de 2011 ``Why Software Is Eating the World'' sostiene que el software se convertirá en la nueva infraestructura sobre la que se apoyan todos los servicios de la industria, y que la capacidad de automatizar la construcción de software será el factor decisivo para la competitividad. Toda la charla de Graham busca llevar esta idea a la práctica en la ingeniería de agentes.
\end{knowledgebox}

\begin{importantbox}{La propuesta de valor del agente}
La promesa central de los agentes de desarrollo de software es convertir un sistema multimodal con capacidad de comprender, modificar, verificar y retroalimentar código en un ejecutor capaz de iterar un producto de manera continua, igual que un desarrollador humano.
\end{importantbox}

\subsection{Resumen de la sección}
\begin{itemize}
    \item El software ya domina la mayoría de las industrias, y la escasez de capacidad de construcción impulsa la necesidad real de agentes.
    \item Graham, al contrastar el Premio Nobel con el ámbito del software, nos recuerda que crear software ya es una de las competencias centrales de la ciencia humana.
\end{itemize}
